% Lebesgue外测度的区间覆盖：E=[0,1]∪{2-2^{-k}:k>=1}；覆盖(-0.1,1.1)与以1.5,1.75,1.875,1.9375为中心、长度0.1,0.05,0.025,0.0125的区间；总长1.2+0.2=1.4。
\begin{tikzpicture}[x=4.4cm,y=1cm,
  cov/.style={draw=FigureBlue,fill=FigureBlueFill,line width=.6pt,rounded corners=1pt},
  note/.style={font=\FigureNoteFont\fontsize{7.5}{9}\selectfont,text=FigureMuted},
  lab/.style={font=\fontsize{7}{8}\selectfont}]
\fill[cov] (-0.1,0.18) rectangle (1.1,0.42);
\node[note] at (0.5,0.68) {长度$1.2$};
\foreach \k/\c/\h in {1/1.5/0.1,2/1.75/0.05,3/1.875/0.025,4/1.9375/0.0125}{\fill[cov] (\c-\h/2,0.18) rectangle (\c+\h/2,0.42);}
\node[note] at (1.72,0.68) {长度$0.1,0.05,0.025,\dots$};
\draw[FigureStroke,line width=.6pt,->] (-0.2,0)--(2.12,0);
\draw[FigureBlue,line width=2.8pt] (0,0)--(1,0);
\foreach \k in {1,...,7}{\fill[FigureBlue] ({2-2^(-\k)},0) circle (1.4pt);}
\foreach \x/\l in {0/0,1/1,1.5/1.5,1.75/1.75,2/2}{\draw[FigureStroke,line width=.5pt] (\x,-.06)--(\x,.06);\node[lab,below] at (\x,-.08) {$\l$};}
\node[note] at (1.05,-0.75) {覆盖总长度$1.2+0.1\,(1+\tfrac12+\tfrac14+\cdots)=1.4$};
\end{tikzpicture}
